\documentclass[12pt]{article}
\usepackage[
  height=9in,      % height of the text block
  width=7in,       % width of the text block
  top=78pt,        % distance of the text block from the top of the page
  headheight=52pt, % height for the header block
  headsep=12pt,    % distance from the header block to the text block
  heightrounded,   % ensure an integer number of lines
  verbose,         % show the values of the parameters in the log file
]{geometry}
\usepackage{array}
\usepackage{multirow}
\usepackage{fancyhdr}
\usepackage{graphicx}
\usepackage{tabularx}
\setlength{\extrarowheight}{5pt}
\usepackage{subcaption}%

% header and footer
\pagestyle{fancy}
\fancyhf{\LARGE\textbf{Learn Basics} \\ \textbf{Chapter Report}}%
\rhead{Subject : \textbf{Math} \\ Class : \textbf{6} \\ Chapter : \textbf{5. Fractions}}%
\lhead{  \includegraphics[width=1.5cm, height=1.5cm]{LearnBasics logo}  }{}%
\lfoot{\textbf{XYZ School}}%
\rfoot{XYZ1118 C6M05 CR##1}%
\cfoot{\thepage}

\renewcommand{\headrulewidth}{2pt}% Default \headrulewidth is 0.4pt
\renewcommand{\footrulewidth}{2pt}% Default \footrulewidth is 0pt

\renewcommand{\arraystretch}{1.2}

\begin{document}

\begin{center}
   \huge{Homework Questions Report}
   \linebreak%
   \underline{\LARGE\textbf{Estimation of Numbers}} 
\end{center}

\vspace{0.5cm}\\%
Concepts Tested in the Question: 

Estimating to the nearest tens by rounding off

Estimating to the nearest hundreds by rounding off

Estimating to the nearest thousands by rounding off

\vspace{1cm}\\

\begin{tabularx}{0.9\textwidth} { 
  | >{\centering\arraybackslash}X 
  | >{\centering\arraybackslash}X | }
 \hline
 Total Students & \textbf{30}\\
 \hline
 No.of Students attended the test  & \textbf{23}\\
\hline
 Absent & \textbf{7} \\
\hline
 Start Time & \textbf{7:00 PM-15th May} \\
\hline
 End Time & \textbf{8:00 PM-16th May} \\
\hline
 Total Questions & \textbf{5-Basic Questions} \\
\hline
\end{tabularx}
\vspace{0.5cm}\\

%%%%%%%%%%%%%%%%%%%%% Two tables side by side %%%%%%%%%%%%%%%%%%%%%%%%%%%%%
\begin{table}
    \begin{subtable}[t]{.5\textwidth}
    \caption{\underline{Top 5 performers :}}
        \raggedright
            \begin{tabularx}{0.9\textwidth} { 
            | >{\centering\arraybackslash}X 
            | >{\centering\arraybackslash}X 
            | >{\centering\arraybackslash}X 
            | >{\centering\arraybackslash}X | }
            \hline
            \textbf{S.No} & \textbf{Name} & \textbf{Score} & \textbf{Time Taken} \\
            \hline
            1 & Arjun & 5 & 5:30min \\
            \hline
            1 & Bhanu & 5 & 5:00min \\
            \hline
            1 & Chandu & 3 & 4:30min \\
            \hline
            1 & Deepu & 4 & 2:30min \\
            \hline
            1 & Varun & 5 & 4:34min \\
            \hline
            \end{tabularx}
    \end{subtable}%
    \begin{subtable}[t]{.4\textwidth}
    \caption{\underline{Absentees :}}
        \raggedleft
        \begin{tabularx}{0.9\textwidth} { 
          | >{\centering\arraybackslash}X 
          | >{\centering\arraybackslash}X | }
         \hline
         \textbf{S.No} & \textbf{Name}\\
         \hline
         1 & Anind \\
         \hline
         2 & Bharath \\
         \hline
         3 & Code \\
         \hline
         4 & Drama \\
         \hline
         5 & Dheeraj \\
         \hline
         6 & Lohith \\
         \hline
         7 & Murali \\
         \hline
        \end{tabularx}
    \end{subtable}
\end{table}
\end{document}